\documentclass{book}

\title{Analisi Matematica T-A}
\author{Michele Mazza}
\date{November 2023}


% Pacchetti di Codifica e Lingua
\usepackage[utf8]{inputenc}
\usepackage[T1]{fontenc}
\usepackage[italian]{babel}

% Pacchetti Grafici e Diagrammi
\usepackage{graphicx}
\usepackage{tikz}
\usepackage{tikz-3dplot}
\usetikzlibrary{angles,quotes,arrows.meta, positioning}
\usepackage{pgfplots}
\pgfplotsset{compat=1.18}

% Pacchetti Matematici
\usepackage{amsmath, amssymb, amsthm}

% Pacchetti per Tabelle e Unità di Misura
\usepackage{geometry}
\usepackage{siunitx}
\usepackage{booktabs}
\usepackage{tabularx}
\usepackage{colortbl}

% Pacchetti per Layout e Formattazione
\usepackage{fancyhdr}
\usepackage{float}
\usepackage{microtype}

% Altri Pacchetti Utili
\usepackage{mdframed}
\usepackage{animate}

% Pacchetto enumitem
\usepackage{enumitem}

% Definizioni dei Colori
\definecolor{lightblue}{HTML}{9cd8fc}
\definecolor{myblue}{HTML}{139FE9}

% Definizione del comando \norm
\newcommand{\norm}[1]{\left\lVert#1\right\rVert}

%\font\nullfont=cmr10 

% Ambienti Personalizzati
\newmdtheoremenv[
    backgroundcolor=orange!10,
    linecolor=myblue,
    linewidth=0pt,
    leftmargin=0pt,
    rightmargin=0pt,
    innerleftmargin=10pt,
    innerrightmargin=10pt,
    innertopmargin=10pt,
    innerbottommargin=10pt,
    nobreak
]{curiosity}{Curiosità}

\newmdtheoremenv[
    font=\normalfont\small,
    linewidth=0pt,
    innerleftmargin=10pt,
    innerrightmargin=10pt,
]{example}{Esempio}

\newmdtheoremenv[
  backgroundcolor=lightblue,
  linecolor=myblue,
  linewidth=1.5pt,
  innertopmargin=10pt,
  innerbottommargin=10pt,
  innerrightmargin=10pt,
  innerleftmargin=10pt,
  nobreak
]{theorem}{Teorema}

\newmdtheoremenv[
  backgroundcolor=gray!10,
  linewidth=0pt,
  innertopmargin=10pt,
  innerbottommargin=10pt,
  innerrightmargin=10pt,
  innerleftmargin=10pt,
  nobreak
]{definition}{Definizione}

\newmdtheoremenv[
  backgroundcolor=gray!10,
  linecolor=myblue,
  linewidth=1.5pt,
  innertopmargin=10pt,
  innerbottommargin=10pt,
  innerrightmargin=10pt,
  innerleftmargin=10pt,
  nobreak
]{proposition}{Proposizione}

% Configurazioni di fancyhdr
\pagestyle{fancy}
\fancyhf{} % Pulisce tutti i campi
\fancyhead[LE,RO]{\thepage} % Numero di pagina
\fancyhead[RE]{\nouppercase{\leftmark}} % Titolo del capitolo
\fancyhead[LO]{\nouppercase{\rightmark}} % Titolo della sezione
\renewcommand{\headrulewidth}{0.5pt}
\renewcommand{\footrulewidth}{0pt}

\geometry{a4paper, left=20mm, right=20mm, top=20mm, bottom=20mm}

% Pacchetto Hyperref (inserito per ultimo)
\usepackage{hyperref}

\begin{document}

\frontmatter
\maketitle
\tableofcontents
\mainmatter

\chapter{Vettori e numeri complessi}

\section{Coordinate polari}
Un punto \(P\) di coordinate \((x , y)\) nel piano cartesiano può essere individuato anche mediante le sue \textbf{coordinate polari} \(r, \theta\).
\begin{definition}[Coordinate polari]
    \((r, \theta)\) sono così definiti,
    \begin{itemize}
        \item \textbf{\(r\)} è la distanza del punto \(P\) dall'origine \(O\).
        \item \textbf{\( \theta \)} è l'angolo in radianti, misurato rispetto al semiasse positivo delle ascisse, che la linea congiungente l'origine \( O \) e il punto \( P \) forma con questo semiasse. L'angolo \( \theta \) è definito modulo \( 2\pi \).\footnote{La notazione "modulo \( 2\pi \)" indica che gli angoli sono misurati in un intervallo da 0 a \( 2\pi \) radianti. Se un angolo supera \( 2\pi \) radianti, viene ridotto sottraendo multipli di \( 2\pi \) fino a rientrare in questo intervallo. Questo è analogo a considerare gli angoli in un cerchio, dove \( 2\pi \) radianti corrispondono a un giro completo di 360 gradi.}
    \end{itemize}
    \textbf{Osservazione :}
    Questo sistema è particolarmente utile per descrivere posizioni e movimenti in scenari circolari o rotazionali.
\end{definition}
\begin{figure}[H]
    \centering
    \begin{tikzpicture}
        \draw[->] (-2.5,0) -- (2.5,0);
        \draw[->] (0,-2.5) -- (0,2.5) node[above] {$y$};
        \draw[->, fill=lightblue] (1.5,0) arc (0:45:1.5) -- (0,0) -- cycle;
        \draw[dashed, myblue] (45:2.9) -- (45:2.9 |- 0,0);
        \draw[dashed, myblue] (45:2.9) -- (45:2.9 -| 0,0);
        \draw[-, thick] (0,0) -- (2,2) node[above right]{$P = ( x , y )$};
        \fill [myblue](2,2) circle (2pt);
        \node at (0,2)[left]{$y$};
        \node at (2,0)[below]{$x$};
        \node at (1.2, 1.2) [above] {$r$};
        \node at(1.4, 0.6)[above right]{$\theta$};
    \end{tikzpicture}
    \caption{Coordinate cartesiane e polari nel piano}
\end{figure}
\(\theta\) viene solitamente scelto nell'intervallo \((-\pi,\pi]\).
\subsection{Passaggio da coordinate polari \texorpdfstring{\((r, \theta)\)}{(r, theta)} a cartesiane \texorpdfstring{\((x, y)\)}{(x, y)}}
\begin{theorem}
    Dato \(P = (r, \theta)\)
    \[x = r \cos{\theta}\]
    \[y = r \sin{\theta}\]
\end{theorem}
\subsection{Passaggio da coordinate cartesiane \texorpdfstring{\((x, y)\)}{(x, y)} a coordinate polari \texorpdfstring{\((r, \theta)\)}{(r, theta)}}
\begin{theorem}
    Dato \(P = (x, y)\) , scelto \(\theta\) nell'intervallo \((-\pi,\pi]\),
    \footnote[1]{Rircordiamo che \(\tan{\theta}=\displaystyle\frac{\sin{\theta}}{\cos{\theta}} = \displaystyle\frac{y}{x} \Rightarrow \theta = \arctan{\displaystyle\frac{y}{x}}\)}
    \begin{align*}
        r = \sqrt{x^2 + y^2}\ ,\  & \ \theta\ =\ \begin{cases}
                                                     \displaystyle \arctan{\frac{y}{x}}       & \text{, se } x > 0            \\[10pt]
                                                     \displaystyle \arctan{\frac{y}{x} + \pi} & \text{, se } x < 0 , y \geq 0 \\[10pt]
                                                     \displaystyle \arctan{\frac{y}{x} - \pi} & \text{, se } x < 0 , y < 0    \\[10pt]
                                                     \displaystyle \frac{\pi}{2}              & \text{, se } x = 0 , y > 0    \\[10pt]
                                                     -\displaystyle \frac{\pi}{2}             & \text{, se } x = 0 , y < 0    \\
                                                 \end{cases}
    \end{align*}
\end{theorem}
\begin{figure}[H]
    \centering
    \begin{tikzpicture}[scale = 0.5]
        \small
        \draw[->] (-2.5,0) -- (2.5,0);
        \draw[->] (0,-2.5) -- (0,2.5);
        \draw[-latex, fill=lightblue, opacity=0.5pt] (0,-1.5) arc (-90:90:1.5) -- (0,0) -- cycle;
        \node at (0,-2.5)[below]{$\displaystyle\arctan{\frac{y}{x}} \text{, se } x > 0  $};
    \end{tikzpicture}
    \begin{tikzpicture}[scale = 0.5]
        \small
        \draw[->] (-2.5,0) -- (2.5,0);
        \draw[->] (0,-2.5) -- (0,2.5);
        \draw[-latex, fill=lightblue, opacity=0.5pt] (0,1.5) arc (90:180:1.5) -- (0,0) -- cycle;
        \node at (0,-2.5)[below]{$\displaystyle\arctan{\frac{y}{x} + \pi} \text{, se } x < 0 , y \geq 0  $};
    \end{tikzpicture}
    \begin{tikzpicture}[scale = 0.5]
        \small
        \draw[->] (-2.5,0) -- (2.5,0);
        \draw[->] (0,-2.5) -- (0,2.5);
        \draw[-latex, fill=lightblue, opacity=0.5pt] (-1.5,0) arc (180:270:1.5) -- (0,0) -- cycle;
        \node at (0,-2.5)[below]{$\displaystyle\arctan{\frac{y}{x} - \pi} \text{, se } x < 0 , y < 0  $};
    \end{tikzpicture}
    \caption{Passaggio da coordinate cartesiane a polari}
\end{figure}
\section{Coordinate cilindriche e sferiche}
\subsection{Coordinate cilindriche}
\begin{figure}[H]
    \def\rotx{70}
    \def\rotz{110}
    \centering
    \tdplotsetmaincoords{\rotx}{\rotz}
    \begin{tikzpicture}[tdplot_main_coords][scale = 0.5]
        \tikzset{every node/.style={font=\small}}
        \draw[dashed,-latex] (0,0,0) -- (6,0,0) node[anchor=north east]{$x$};
        \draw[dashed,-latex] (0,0,0) -- (0,6,0) node[anchor=north west]{$y$};
        \draw[dashed,-latex] (0,0,0) -- (0,0,6) node[anchor=south]{$z$};
        \draw [thick](0,0,0) circle (3);
        \draw [thick](0,0,4) circle (3);
        \filldraw[fill=gray, nearly transparent] (-4,-4,4) -- (4,-4,4) --  (4,5,4) -- (-4,5,4) -- (-4,-4,4);
        \filldraw[fill=gray, nearly transparent] (0,0,4) -- (5.2,6,4) --  (5.2,6,0) -- (0,0,0) -- (0,0,4);
        \draw [thick, myblue] (0,0,0) -- (2,2.25,4);
        \node at (2,2.25,4.5)[color=myblue,above] {$P = ( r, \varphi, z )$};
        \filldraw [color=myblue](2,2.25,4) circle (0.075cm);

        % manual edges
        \draw [dotted](1.9,-2.34,0) -- (1.9,-2.34,4) node[midway, left]{};
        \draw [dotted](-1.9,2.35,0) -- (-1.9,2.35,4);
        % automatic edges !
        \pgfcoordinate{edge1_top}{ \pgfpointcylindrical{\rotz}{3}{4} };
        \pgfcoordinate{edge1_bottom}{ \pgfpointcylindrical{\rotz}{3}{0} };
        \draw[thick] (edge1_top) -- (edge1_bottom);
        \pgfcoordinate{edge2_top}{ \pgfpointcylindrical{\rotz+180}{3}{4} };
        \pgfcoordinate{edge2_bottom}{ \pgfpointcylindrical{\rotz+180}{3}{0} };
        \draw[thick] (edge2_top) -- (edge2_bottom);
    \end{tikzpicture}
    \caption{Coordinate cilindriche}
\end{figure}
\begin{theorem}
    Dato \(P = (r, \varphi, z)\)
    \[x = r \cos{\varphi}\]
    \[y = r \sin{\varphi}\]
    \[z = z\]
\end{theorem}
\subsection{Coordinate sferiche}
\begin{figure}[H]
    \centering
    \begin{tikzpicture}
        \tikzstyle{every node}=[font=\small]
        % Definisci il raggio
        \def\r{3}
        % Vettore Bloch
        \draw[thick, myblue] (0, 0) node[circle, fill, inner sep=1] (orig) {} -- (\r/3, \r/2) node[circle, fill, inner sep=1,] (a) {};
        \node at (\r/3, \r/2)[myblue,above right] {$P = ( r, \varphi, \theta )$};
        \node at (\r/4, \r/3) [color = myblue, above left] {$r$};
        \draw[dashed] (orig) -- (\r/3, -\r/5) node (phi) {} -- (a);
        % Sfera con ombreggiatura
        \shade[ball color = gray, opacity = 0.2] (orig) circle (\r);
        \draw[dashed] (orig) ellipse (\r{} and \r/3);
        % Assi
        \draw[->] (orig) -- ++(-\r/5, -\r/3) node[below] (x) {$x$};
        \draw[->] (orig) -- ++(\r, 0) node[right] (y) {$y$};
        \draw[->] (orig) -- ++(0, \r) node[above] (z) {$z$};
        % Angoli
        \pic [draw=myblue, text=myblue, ->, "$\theta$"] {angle = x--orig--phi};
        \pic [draw=myblue, text=myblue, <-, "$\varphi$", angle eccentricity=1.4] {angle = a--orig--z};
    \end{tikzpicture}
    \caption{Coordinate sferiche nella sfera di Bloch}
\end{figure}
\begin{theorem}
    Dato \(P = (r, \varphi, \theta)\)
    \[x = r \sin{\varphi} \cos{\theta}\]
    \[y = r \sin{\varphi} \sin{\theta}\]
    \[z = r \cos{\varphi}\]
\end{theorem}
\section{Vettori nel piano e nello spazio}
\begin{definition}[Vettore]
    Un \textbf{vettore} \(\vec{v}=( v_1, v_2, ..., v_n)\) è un segmento orientato, cioè un segmento che ha un \textbf{modulo} (lunghezza) e una \textbf{direzione} (retta su cui giace) e un \textbf{verso} (orientamento).
    \begin{itemize}
        \item \textbf{Modulo:} è la lunghezza del vettore.
        \item \textbf{Direzione:} è la retta su cui giace il vettore.
        \item \textbf{Verso:} è il senso di percorrenza della retta su cui giace.
    \end{itemize}
\end{definition}
\begin{figure}[H]
    \centering
    \begin{tikzpicture}
        \draw[->] (-2.5,0) -- (2.5,0);
        \draw[->] (0,-2.5) -- (0,2.5);
        \draw[-latex, ultra thick, myblue] (0,0) -- (2,2);
        \node at (2,2)[above right]{$P = ( v_1 , v_2 )$};
        \node at (1,1)[above left]{$\vec{v}$};
        \node at (0,0)[below left]{$O$};
        \node at (0,2)[left]{$y$};
        \node at (2,0)[below]{$x$};
    \end{tikzpicture}
    \begin{tikzpicture}[scale=1]
        \draw[->] (0,0,0) -- (3,0,0);
        \draw[->] (0,0,0) -- (0,3,0);
        \draw[->] (0,0,0) -- (0,0,3);
        \draw[-latex, ultra thick, myblue] (0,0,0) -- (2,2,1);
        \node at (0,0,0) [left] {$O$};
        \node at (0,3,0) [left] {$y$};
        \node at (3,0,0) [below] {$x$};
        \node at (0,0,3) [above left] {$z$};
        \node at (2,2,1) [above right] {$P = (w_1, w_2, w_3)$};
        \node at (1,1,0.5) [below] {$\vec{w}$};
    \end{tikzpicture}
    \caption{Vettori nel piano e nello spazio}
    \label{fig:vettori}
\end{figure}
\subsection{Proprietà dei vettori}
\begin{definition}[Norma euclidea]
    La \textbf{norma} di un vettore \(\vec{v}=(v_1, v_2, \ldots, v_n)\) è il numero reale non negativo \(\norm{\vec{v}}\) definito come
    \[
        \norm{\vec{v}} = \sqrt{\sum_{i=1}^{n} v_i^2} = \sqrt{v_1^2 + v_2^2 + \ldots + v_n^2}
    \]

    \noindent Osservazioni e Proprietà:
    \begin{itemize}
        \item La norma euclidea misura la "lunghezza" o la "distanza" di un vettore dall'origine nello spazio n-dimensionale.
        \item Soddisfa la disuguaglianza triangolare: \(\norm{\vec{v} + \vec{w}} \leq \norm{\vec{v}} + \norm{\vec{w}}\) per ogni coppia di vettori \(\vec{v}\) e \(\vec{w}\).
        \item \(\norm{\vec{v}} = 0\) se e solo se \(\vec{v}\) è il vettore nullo.
        \item È una funzione omogenea: \(\norm{k\vec{v}} = |k| \norm{\vec{v}}\) per ogni scalare \(k\).
        \item La norma euclidea deriva direttamente dal prodotto scalare: \(\norm{\vec{v}} = \sqrt{\vec{v} \cdot \vec{v}}\).
    \end{itemize}
\end{definition}

\begin{definition}[Versore]
    Un \textbf{versore} è un vettore di norma 1.

    \noindent Osservazioni e Proprietà:
    \begin{itemize}
        \item Un versore è un vettore unitario, cioè ha lunghezza (o norma) pari a 1.
        \item È spesso utilizzato per indicare una direzione nello spazio, poiché la sua lunghezza unitaria rende facile considerare solo la direzione senza influenze della magnitudine.
        \item Ogni vettore non nullo può essere trasformato in un versore dividendo il vettore per la sua norma: se \(\vec{v}\) è un vettore, allora \(\frac{\vec{v}}{\norm{\vec{v}}}\) è un versore.
        \item I versori sono particolarmente utili nella definizione di sistemi di coordinate e nella descrizione di basi ortonormali.
    \end{itemize}
\end{definition}

\subsection{Operazioni tra vettori}
\subsubsection{Somma tra vettori}
\begin{theorem}[Somma tra vettori]
    Dati due vettori \(\vec{v}=(v_1, v_2, \ldots, v_n)\) e \(\vec{w}=(w_1, w_2, \ldots, w_n)\), la loro \textbf{somma} è il vettore
    \[
        \vec{v}+\vec{w}=(v_1+w_1, v_2+w_2, \ldots, v_n+w_n)
    \]

    \noindent Osservazioni e Proprietà:
    \begin{itemize}
        \item La somma di vettori segue la regola del parallelogramma: la somma di due vettori può essere rappresentata geometricamente come il completamento di un parallelogramma i cui lati sono i due vettori.
        \item La somma è commutativa, ovvero \(\vec{v} + \vec{w} = \vec{w} + \vec{v}\).
        \item La somma è anche associativa, cioè \((\vec{v} + \vec{w}) + \vec{u} = \vec{v} + (\vec{w} + \vec{u})\).
        \item Aggiungendo lo zero vettoriale, \(\vec{0}\), a qualsiasi vettore \(\vec{v}\), il risultato sarà \(\vec{v}\) stesso: \(\vec{v} + \vec{0} = \vec{v}\).
    \end{itemize}
\end{theorem}

\begin{figure}[H]
    \centering
    \begin{animateinline}[controls,loop]{60}
        \multiframe{10}{rW=-1.0+0.3}{%
            \begin{tikzpicture}
                \draw[->] (-2.5,0) -- (2.5,0);
                \draw[->] (0,-2.5) -- (0,2.5);
                \draw[dashed] (-2.5,-2.5) -- (2.5,2.5);
                \draw[dashed] (-2.5,1) -- (3,1);
                \draw[dashed] (-2.5+\rW,-2.5) -- (2.5+\rW,2.5);
                \draw[-latex, ultra thick, myblue] (0,0) -- (1,1); % Vettore v
                \draw[-latex, ultra thick, myblue] (0,0) -- (\rW,0); % Vettore w con componente orizzontale variabile
                \draw[-latex, ultra thick, myblue] (0,0) -- ({1+\rW},1); % Vettore v+w
                \node at (0,0)[below left]{$O$};
                \node at (0.5,0.5)[above left]{$\vec{v}$};
                \node at (\rW/2,0)[below]{$\vec{w}$};
                \node at ({(1+\rW)/2},0.5)[right]{$\vec{v}+\vec{w}$};
                \node at (2.5,-2.5,0)[below]{$\vec{v}=(1,1)$ $\vec{w}=(\rW,0)$};
            \end{tikzpicture}%
        }
    \end{animateinline}
    \begin{tikzpicture}
        font=\small
        \draw[->] (-2.5,0) -- (2.5,0);
        \draw[->] (0,-2.5) -- (0,2.5);
        \draw[dashed] (-2.5,-2.5) -- (2.5,2.5);
        \draw[dashed] (-2.5,1) -- (3,1);
        \draw[dashed] (-0.5,-2.5) -- (4.5,2.5);
        \draw[-latex, ultra thick, myblue] (0,0) -- (1,1);
        \draw[-latex, ultra thick, myblue] (0,0) -- (2,0);
        \draw[-latex, ultra thick, myblue] (0,0) -- (3,1);
        \node at (0,0)[below left]{$O$};
        \node at (0.5,0.5)[above left]{$\vec{v}$};
        \node at (1,0)[below]{$\vec{w}$};
        \node at (3,1)[right]{$\vec{v}+\vec{w}$};
    \end{tikzpicture}
    \caption{Somma tra vettori}
    \label{fig:somma}
\end{figure}

\subsubsection{Prodotto per uno scalare}
\begin{theorem}[Prodotto per uno scalare]
    Dato un vettore \(\vec{v}=(v_1, v_2, \ldots, v_n)\) e uno scalare \(\lambda\), il loro \textbf{prodotto} è il vettore
    \[
        \lambda\vec{v}=(\lambda v_1, \lambda v_2, \ldots, \lambda v_n)
    \]

    \noindent Osservazioni e Proprietà:
    \begin{itemize}
        \item Il prodotto di un vettore per uno scalare cambia la lunghezza del vettore senza alterare la sua direzione, a meno che lo scalare non sia negativo, nel qual caso il vettore viene anche invertito.
        \item Questa operazione è nota come \textit{moltiplicazione scalare} e è una delle operazioni fondamentali nell'algebra lineare.
        \item La moltiplicazione scalare è distributiva rispetto alla somma di vettori: \(\lambda(\vec{v} + \vec{w}) = \lambda\vec{v} + \lambda\vec{w}\).
        \item È anche distributiva rispetto alla somma di scalari: \((\lambda + \mu)\vec{v} = \lambda\vec{v} + \mu\vec{v}\) per ogni scalare \(\mu\).
        \item Moltiplicare un vettore per 1 lascia il vettore invariato, mentre moltiplicarlo per 0 dà il vettore nullo: \(1\vec{v} = \vec{v}\) e \(0\vec{v} = \vec{0}\).
    \end{itemize}
\end{theorem}
\begin{figure}[H]
    \centering
    \begin{animateinline}[controls,loop]{60}
        \multiframe{10}{rLambda=-1+0.4}{%
            \begin{tikzpicture}
                \draw[->] (-2.5,0) -- (2.5,0);
                \draw[->] (0,-2.5) -- (0,2.5);
                \draw[dashed] (-2.5,-2.5) -- (2.5,2.5);
                \draw[-latex, ultra thick, myblue] (0,0) -- (0.5,0.5); % Vettore v
                \draw[-latex, ultra thick, myblue] (0,0) -- (\rLambda*0.5,\rLambda*0.5); % Vettore lambda*v
                \node at (0,0)[below left]{$O$};
                \node at (0.5,0.5)[above left]{$\vec{v}$};
                \node at (\rLambda*0.5,\rLambda*0.5)[below right]{$\rLambda\vec{v}$};
            \end{tikzpicture}%
        }
    \end{animateinline}
    \begin{tikzpicture}
        font=\small
        \draw[->] (-2.5,0) -- (2.5,0);
        \draw[->] (0,-2.5) -- (0,2.5);
        \draw[dashed] (-2.5,-2.5) -- (2.5,2.5);
        \draw[-latex, ultra thick, myblue] (0,0) -- (1.5,1.5);
        \draw[-latex, ultra thick, myblue] (0,0) -- (0.5,0.5);
        \node at (0,0)[below left]{$O$};
        \node at (0.5,0.5)[above left]{$\vec{v}$};
        \node at (1.5,1.5)[below right]{$3\vec{v}$};
    \end{tikzpicture}
    \caption{Prodotto per uno scalare}
    \label{fig:scalare}
\end{figure}
\subsubsection{Prodotto scalare}
\noindent\begin{minipage}{\textwidth}
    \begin{theorem}[Prodotto scalare]
        Dati due vettori \(\mathbf{v}\) e \(\mathbf{w}\), il loro \textbf{prodotto scalare} è lo scalare (dunque un numero non un vettore) definito come:
        \[
            \mathbf{v} \cdot \mathbf{w} = \sum_{i=1}^{n} v_i w_i = v_1w_1 + v_2w_2 + \cdots + v_nw_n
        \]
        In termini geometrici, può essere espresso anche come:
        \[
            \mathbf{v} \cdot \mathbf{w} = \|\mathbf{v}\| \|\mathbf{w}\| \cos(\theta)
        \]
        dove \(\theta\) è l'angolo tra i vettori.
        \noindent Osservazioni e Proprietà:
        \begin{itemize}
            \item Il prodotto scalare è commutativo: \(\mathbf{v} \cdot \mathbf{w} = \mathbf{w} \cdot \mathbf{v}\).
            \item Il prodotto scalare è distributivo rispetto alla somma vettoriale: \(\mathbf{v} \cdot (\mathbf{w} + \mathbf{u}) = \mathbf{v} \cdot \mathbf{w} + \mathbf{v} \cdot \mathbf{u}\).
            \item Il prodotto scalare di un vettore con se stesso è uguale alla norma al quadrato di quel vettore: \(\mathbf{v} \cdot \mathbf{v} = \|\mathbf{v}\|^2\).
            \item Se il prodotto scalare di due vettori è zero, allora i vettori sono perpendicolari (ortogonali) tra loro.
            \item Il prodotto scalare misura l'allineamento tra due vettori e può essere usato per calcolare la proiezione di un vettore su un altro. In fisica, è spesso utilizzato per calcolare il lavoro svolto da una forza lungo un percorso.
        \end{itemize}
    \end{theorem}
\end{minipage}
\begin{figure}[H]
    \centering
    \begin{animateinline}[controls,loop]{20}
        \multiframe{10}{rwx=-1+0.3}{%
            \begin{tikzpicture}
                font=\small
                \def\rBound{3}
                \def\vX{2}
                \def\vY{2}
                \def\wY{0.5}
                \def\projWvX{\vX * (\vX*\rwx + \vY*\wY) / (\vX*\vX + \vY*\vY)}
                \def\projWvY{\vY * (\vX*\rwx + \vY*\wY) / (\vX*\vX + \vY*\vY)}
                \draw[->] (-\rBound+2,0) -- (\rBound+3.5,0);
                \draw[->] (0,-\rBound) -- (0,\rBound);
                \draw[dashed] (-\rBound+2,-\rBound+2) -- (\rBound,\rBound);
                \draw[-latex, ultra thick, myblue] (0,0) -- (\vX,\vY);
                \draw[-latex, ultra thick, myblue] (0,0) -- (\rwx,\wY);
                \node at ({\vX/2},{\vY/2})[above left]{$\vec{v}$};
                \node at ({\rwx/2},{\wY/2})[below right]{$\vec{w}$};
                \draw[dashed] (\rwx,\wY) -- ({\projWvX},{\projWvY});
                \draw[-latex, ultra thick] (0,0) -- ({\projWvX},{\projWvY});
                \node at ({\projWvX/2},{\projWvY/2})[above left]{$\vec{w_x}$};
                \fill[myblue] ({(\vX*\rwx)+(\vY*\wY)},0) circle (2pt);
                \node at ({(\vX*\rwx)+(\vY*\wY)},0)[below]{$\vec{v} \cdot \vec{w}$};
            \end{tikzpicture}
        }
    \end{animateinline}
    \begin{tikzpicture}
        font=\small
        \def\rBound{3}
        \def\vX{2}
        \def\vY{2}
        \def\rwx{2.5}
        \def\wY{0.5}
        \def\projWvX{\vX * (\vX*\rwx + \vY*\wY) / (\vX*\vX + \vY*\vY)}
        \def\projWvY{\vY * (\vX*\rwx + \vY*\wY) / (\vX*\vX + \vY*\vY)}
        \draw[->] (-\rBound+2,0) -- (\rBound+3.5,0);
        \draw[->] (0,-\rBound) -- (0,\rBound);
        \draw[dashed] (-\rBound+2,-\rBound+2) -- (\rBound,\rBound);
        \draw[-latex, ultra thick, myblue] (0,0) -- (\vX,\vY);
        \draw[-latex, ultra thick, myblue] (0,0) -- (\rwx,\wY);
        \node at ({\vX/2},{\vY/2})[above left]{$\vec{v}$};
        \node at ({\rwx/2},{\wY/2})[below right]{$\vec{w}$};
        \draw[dashed] (\rwx,\wY) -- ({\projWvX},{\projWvY});
        \draw[-latex, ultra thick] (0,0) -- ({\projWvX},{\projWvY});
        \node at ({\projWvX/2},{\projWvY/2})[above left]{$\vec{w_x}$};
        \fill[myblue] ({(\vX*\rwx)+(\vY*\wY)},0) circle (2pt);
        \node at ({(\vX*\rwx)+(\vY*\wY)},0)[below]{$\vec{v} \cdot \vec{w}$};
    \end{tikzpicture}
    \caption{Prodotto scalare}
\end{figure}
\subsubsection{Prodotto vettoriale}
\noindent\begin{minipage}{\textwidth}
    \begin{theorem}[Prodotto vettoriale]
        Dati due vettori \(\mathbf{u}\) e \(\mathbf{v}\) in \(\mathbb{R}^3\), il loro \textbf{prodotto vettoriale} è un vettore \(\mathbf{w}\), definito come:
        \[
            \mathbf{w} = \mathbf{u} \times \mathbf{v} = \det
            \begin{vmatrix}
                \mathbf{i} & \mathbf{j} & \mathbf{k} \\
                u_1        & u_2        & u_3        \\
                v_1        & v_2        & v_3        \\
            \end{vmatrix}
        \]
        dove \(\mathbf{i}\), \(\mathbf{j}\), e \(\mathbf{k}\) sono i vettori unitari lungo gli assi x, y, e z, rispettivamente. Il prodotto vettoriale può essere calcolato come il determinante di questa matrice \(3 \times 3\).

        In termini geometrici, il prodotto vettoriale \(\mathbf{u} \times \mathbf{v}\) ha le seguenti caratteristiche:
        \begin{itemize}
            \item È perpendicolare sia a \(\mathbf{u}\) che a \(\mathbf{v}\).
            \item La sua norma è data da \(\|\mathbf{u} \times \mathbf{v}\| = \|\mathbf{u}\| \|\mathbf{v}\| \sin(\theta)\), dove \(\theta\) è l'angolo tra \(\mathbf{u}\) e \(\mathbf{v}\).
            \item Segue la regola della mano destra per determinare la sua direzione.
        \end{itemize}
        \noindent Osservazioni e Proprietà:
        \begin{itemize}
            \item Il prodotto vettoriale è anticommutativo: \(\mathbf{u} \times \mathbf{v} = -(\mathbf{v} \times \mathbf{u})\).
            \item È distributivo rispetto alla somma vettoriale: \(\mathbf{u} \times (\mathbf{v} + \mathbf{w}) = \mathbf{u} \times \mathbf{v} + \mathbf{u} \times \mathbf{w}\) per ogni \(\mathbf{w} \in \mathbb{R}^3\).
            \item Se \(\mathbf{u}\) e \(\mathbf{v}\) sono paralleli (o uno dei due è il vettore nullo), il loro prodotto vettoriale è il vettore nullo.
            \item Il prodotto vettoriale può essere usato per calcolare l'area di un parallelogramma formato dai vettori \(\mathbf{u}\) e \(\mathbf{v}\).
        \end{itemize}
    \end{theorem}
\end{minipage}
\begin{figure}[H]
    \centering
    \begin{tikzpicture}[font=\small]
        \def\rBound{3}
        \def\uX{2}
        \def\uY{1}
        \def\vX{1}
        \def\vY{2}
        \def\wZ{2}
        \draw[->] (-\rBound,0,0) -- (\rBound,0,0); % Asse X
        \draw[->] (0,-\rBound,0) -- (0,\rBound,0); % Asse Y
        \draw[->] (0,0,-\rBound) -- (0,0,\rBound); % Asse Z
        \draw[-{Latex[length=3mm,width=2mm]}, ultra thick, myblue] (0,0,0) -- (\uX,\uY,0); % vettore u
        \draw[-{Latex[length=3mm,width=2mm]}, ultra thick, myblue] (0,0,0) -- (\vX,\vY,0); % vettore v
        \draw[-{Latex[length=3mm,width=2mm]}, ultra thick, myblue] (0,0,0) -- (0,0,\wZ); % vettore u x v
        \node at (\uX,\uY,0) [right] {$\mathbf{u}$};
        \node at (\vX,\vY,0) [above] {$\mathbf{v}$};
        \node at (0,0,\wZ) [left] {$\mathbf{u} \times \mathbf{v}$};
    \end{tikzpicture}
    \caption{Prodotto vettoriale}
\end{figure}

\section{Numeri complessi}
L'invenzione dei numeri reali nasce dalla necessità di risolvere le equazioni di secondo grado con discriminante negativo, ovvero
\[ax^2 + bx + c = 0 \text{ con } \Delta < 0\]
come ad esempio
\[x^2 + 1 = 0\]
\noindent\begin{minipage}{\textwidth}
    \begin{definition}[Numero complesso]
        Un \textbf{numero complesso} \( z \) è una coppia ordinata di numeri reali \( (x, y) \), dove
        \begin{enumerate}
            \item \( x \) è la \textbf{parte reale}, \( \Re(z) = x \)
            \item \( y \) è la \textbf{parte immaginaria}, \( \Im(z) = y \).
        \end{enumerate}
        \noindent Osservazioni e Proprietà:
        \begin{itemize}
            \item  L'insieme dei numeri complessi, indicato con \( \mathbb{C} \), è rappresentabile come \( \mathbb{R}^2 \).\footnote{Questa corrispondenza tra \( \mathbb{C} \) e \( \mathbb{R}^2 \) si basa sulla rappresentazione geometrica e algebrica dei numeri complessi. Ogni numero complesso può essere espresso come \( z = x + yi \), dove \( x \) e \( y \) sono numeri reali e \( i \) è l'unità immaginaria. Geometricamente, questo si traduce nel piano complesso, dove ogni punto \( (x, y) \) rappresenta un numero complesso unico, con \( x \) e \( y \) che corrispondono rispettivamente alle parti reale e immaginaria di \( z \). Pertanto, \( \mathbb{C} \) funziona come uno spazio bidimensionale analogo a \( \mathbb{R}^2 \).}
            \item   I numeri della forma \( (x, 0) \) appartengono all'insieme dei numeri reali \( \mathbb{R} \subseteq \mathbb{C} \), mentre quelli della forma \( (0, y) \) sono \textbf{immaginari puri}.
            \item Due numeri complessi \( z_1 = (x_1, y_1) \) e \( z_2 = (x_2, y_2) \) sono uguali se e solo se \( x_1 = x_2 \) e \( y_1 = y_2 \).
        \end{itemize}
    \end{definition}
\end{minipage}
\begin{definition}[Unità immaginaria]
    \textbf{L'unità immaginaria} \(i\) è il numero complesso, \((0, 1)\), esso gode della proprietà
    \[ i^2 = (0,1)\cdot(0,1) = (-1,0) = -1, \]
    \[ \pm i = \sqrt{-1} \]
\end{definition}



\subsection{Notazione in coppie ordinate}
La notazione in coppie ordinate è utile principalmente per enfatizzare l'aspetti geometrico dei numeri complessi. In questa notazione, un numero complesso \( z \) è rappresentato come una coppia ordinata di numeri reali \( (x, y) \), dove \( x \) è la parte reale e \( y \) è la parte immaginaria. In questa notazione, la somma e il prodotto di due numeri complessi \( z_1 = (x_1, y_1) \) e \( z_2 = (x_2, y_2) \) sono definiti come
\begin{theorem}[Somma tra numeri complessi]
    Dati due numeri complessi \( z_1 = (x_1, y_1) \) e \( z_2 = (x_2, y_2) \), la loro \textbf{somma} è il numero complesso
    \[ z_1 + z_2 = (x_1 + x_2, y_1 + y_2). \]
\end{theorem}
\begin{theorem}[Prodotto tra numeri complessi]
    Dati due numeri complessi \( z_1 = (x_1, y_1) \) e \( z_2 = (x_2, y_2) \), il loro \textbf{prodotto} è il numero complesso
    \[ z_1 \cdot z_2 = (x_1 x_2 - y_1 y_2, x_1 y_2 + x_2 y_1). \]
\end{theorem}

\subsection{Notazione algebrica o cartesiana}
La notazione algebrica è la forma più utilizzata per rappresentare i numeri complessi. In questa notazione, un numero complesso \( z \) è rappresentato come una somma di due numeri reali, uno dei quali è moltiplicato per l'unità immaginaria \( i \). In questa notazione, la somma e il prodotto di due numeri complessi \( z_1 = x_1 + i y_1 \) e \( z_2 = x_2 + i y_2 \) sono definiti come.

\noindent\begin{minipage}{\textwidth}
    \begin{definition}[Forma cartesiana o algebrica]
        Un numero complesso \( z \) può essere scritto nella \textbf{forma cartesiana} o \textbf{algebrica} come
        \[ z = (x, y) = x + iy, \footnote{Tale equzione si ricava da \[(x,y) = (x, 0) + (0, 1) (y, 0) = (x, 0) + ( 0 \cdot y - 1 \cdot 0,\ 0 \cdot 0 + y \cdot 1) = (x, 0) + (0, y) = (x, y) \Rightarrow x +iy\]} \]
        dove \( x \) è la parte reale e \( y \) è la parte immaginaria.
    \end{definition}
\end{minipage}

\begin{theorem}[Somma tra numeri complessi]
    Dati due numeri complessi \( z_1 = x_1 + i y_1 \) e \( z_2 = x_2 + i y_2 \), la loro \textbf{somma} è il numero complesso
    \[ z_1 + z_2 = (x_1 + x_2) + i (y_1 + y_2). \]
\end{theorem}
\begin{theorem}
    Dati due numeri complessi \( z_1 = x_1 + i y_1 \) e \( z_2 = x_2 + i y_2 \), il loro \textbf{prodotto} è il numero complesso
    \[ z_1 \cdot z_2 = (x_1 x_2 - y_1 y_2) + i (x_1 y_2 + x_2 y_1). \]
\end{theorem}
\subsection{Proprietà dei numeri complessi}
\begin{theorem}[Proprietà dei numeri complessi]
    Dati tre numeri complessi \( z_1 = x_1 + i y_1 \), \( z_2 = x_2 + i y_2 \) e \( z_3 = x_3 + i y_3 \), valgono le seguenti proprietà:
    \begin{itemize}
        \item Somma: \( z_1 + z_2 = z_2 + z_1 \)
        \item Prodotto: \( z_1 \cdot z_2 = z_2 \cdot z_1 \)
        \item Associatività della somma: \( (z_1 + z_2) + z_3 = z_1 + (z_2 + z_3) \)
        \item Associatività del prodotto: \( (z_1 \cdot z_2) \cdot z_3 = z_1 \cdot (z_2 \cdot z_3) \)
        \item Distributività del prodotto rispetto alla somma: \( z_1 \cdot (z_2 + z_3) = z_1 \cdot z_2 + z_1 \cdot z_3 \)
        \item Elemento neutro della somma: \( z + 0 = z \)
        \item Elemento neutro del prodotto: \( z \cdot 1 = z \)
        \item Elemento opposto della somma: \( z + (-z) = 0 \)
        \item Elemento inverso del prodotto: \( z \cdot \frac{1}{z} = 1 \)
    \end{itemize}
\end{theorem}
\subsubsection{Coniugato di un numero complesso}
\begin{definition}[Coniugato di un numero complesso]
    Il \textbf{coniugato} di un numero complesso \( z = x + iy \) è il numero complesso
    \[ \overline{z} = x - iy. \]
\end{definition}
\subsubsection{Modulo di un numero complesso}
\noindent\begin{minipage}{\textwidth}
    \begin{definition}[Modulo di un numero complesso]
        Il \textbf{modulo} di un numero complesso \( z = x + iy \) è il numero reale
        \[ |z| = \sqrt{x^2 + y^2}. \]
    \end{definition}
\end{minipage}
\subsection{Coordinate cartesiane}
È naturale associare al numero \( z = (x, y) = x + iy \) il punto del piano cartesiano
di coordinate \( x \) e \( y \) (si veda la Figura \ref{fig:cartesiane}). Il numero \( z \) può anche essere
pensato come il vettore applicato nell'origine e avente tale punto come estremo.
\begin{figure}[H]
    \centering
    \begin{tikzpicture}
        \draw[->] (-2.5,0) -- (2.5,0);
        \draw[->] (0,-2.5) -- (0,2.5);
        \draw[dashed, myblue] (2,2) -- (2,0);
        \draw[dashed, myblue] (2,2) -- (0,2);
        \draw[dashed, myblue] (2,-2) -- (2,0);
        \draw[dashed, myblue] (2,-2) -- (0,-2);
        \node at(2,2)[above right]{$z = x + iy$};
        \node at(2,-2)[below right]{$\overline{z} = x - iy$};
        \fill [myblue](2,2) circle (2pt);
        \fill [myblue](2,-2) circle (2pt);
        \node at (0, 2.5)[above right]{$\Im(z)$};
        \node at (2.5, 0)[above right]{$\Re(z)$};
        \node at (0,2)[left]{$y$};
        \node at (2,0)[below]{$x$};
        \node at (0,0)[below left]{$O$};
    \end{tikzpicture}
    \caption{Coordinate cartesiane di un numero complesso e il suo coniugato}
    \label{fig:cartesiane}
\end{figure}
\subsection{Notazione trigonometrica o polare}
\begin{definition}[Notazione trigonometrica o polare]
    Un numero complesso \( z \) può essere scritto nella \textbf{forma trigonometrica} come
    \[ z = r(\cos(\theta) + i \sin(\theta)), \]
    dove \( r = |z| \) e \( \theta \) è l'angolo che \( z \) forma con l'asse reale positivo.

    \noindent Osservazioni e Proprietà:
    \begin{itemize}
        \item \(\theta\) è detto \textbf{argomento} di \(z\), e si indica con \(\theta = \arg(z)\). \\
        \item \(-\pi < \theta \leq pi \) è \textbf{l'argomento principale} di \(z\), e si indica con \(\theta = Arg(z)\).
    \end{itemize}
    \footnote{La rappresentazione polare risulta utile per esprimere in maniera semplice il prodotto di due numeri complessi.}
\end{definition}
\begin{figure}[H]
    \centering
    \begin{tikzpicture}
        \draw[->] (-2.5,0) -- (2.5,0);
        \draw[->] (0,-2.5) -- (0,2.5);
        \draw[dashed, myblue] (2,2) -- (2,0);
        \draw[dashed, myblue] (2,2) -- (0,2);
        \draw[thick] (0,0) -- (2,2);
        \draw[-latex] (1,0) arc (0:45:1);
        \node at (1,1)[below right] {$\theta$};
        \node at(1,1)[above left]{$r$};
        \node at(2,2)[above right]{$z = r(\cos(\theta) + i \sin(\theta))$};
        \fill [myblue](2,2) circle (2pt);
        \node at (0, 2.5)[above right]{$\Im(z)$};
        \node at (2.5, 0)[above right]{$\Re(z)$};
        \node at (0,2)[left]{$y$};
        \node at (2,0)[below]{$x$};
        \node at (0,0)[below left]{$O$};
    \end{tikzpicture}
    \caption{Coordinate polari di un numero complesso}
    \label{fig:polari}
\end{figure}
\subsubsection{Prodotto di numeri complessi in forma trigonometrica}
\noindent\begin{minipage}{\textwidth}
    \begin{theorem}[Prodotto di numeri complessi in forma trigonometrica]
        Dati due numeri complessi \( z_1 = r_1(\cos\theta_1 + i \sin\theta_1) \) e \( z_2 = r_2(\cos\theta_2 + i \sin\theta_2) \), il loro \textbf{prodotto} è il numero complesso
        \[ z_1 \cdot z_2 = r_1 r_2 (\cos(\theta_1 + \theta_2) + i \sin(\theta_1 + \theta_2)). \]
    \end{theorem}
\end{minipage}
\begin{proposition}
    Vale la seguente proprietà:
    \[\text{arg}\left(z_1 z_2\right) = \text{arg}\left(z_1\right) + \text{arg}\left(z_2\right)\]
    che evidenzia la natura rotazionale del prodotto di numeri complessi.
\end{proposition}
\begin{curiosity}[Moltiplicare per \(i\)]
    \'E giunto il momento di parlare di un aspetto importante del significato di \(i\). Consideriamo il numero complesso \(z\) nelle tre forme che abbiamo incontrato finora:
    \begin{align*}
        z & = x + iy \text{ (forma algebrica)}                                \\
        z & = (x, y) \text{ (coppie ordinate)}                                \\
        z & = r(\cos(\theta) + i \sin(\theta)) \text{ (forma trigonometrica)}
    \end{align*}
    ecco come si comporta la moltiplicazione per \(i\) in ciascuna di queste forme:
    \begin{align*}
        iz & = i(x + iy) = ix + i^2y = -y + ix \text{ (forma algebrica)}                                               \\
        iz & = (0,1)(x,y) = (-y, x) \text{ (coppie ordinate)}                                                          \\
        iz & = 1 \cdot r(\cos(\frac{\pi}{2} + \theta) + i \sin(\frac{\pi}{2} + \theta)) \text{ (forma trigonometrica)}
    \end{align*}
    Graficamente questi risultati si possono interpretare come rotazioni di \(z\) di \(90^\circ\) in senso antiorario del piano di Gauss.
    \begin{figure}[H]
        \centering
        \begin{tikzpicture}
            \draw[->] (-2.5,0) -- (2.5,0);
            \draw[->] (0,-2.5) -- (0,2.5);
            \draw[thick] (0,0) -- (2,1);
            \draw[thick] (0,0) -- (-1,2);
            \draw[-latex] (1,0) arc (0:30:1);
            \draw[-latex] (1,0.5) arc (30:120:1);
            \node at (1,0.5)[below right] {$\theta$};
            \node at(2,1)[above right]{$z = r(\cos(\theta) + i \sin(\theta))$};
            \node at(-1,2)[above left]{$iz = r(\cos(\frac{\pi}{2} + \theta) + i \sin(\frac{\pi}{2} + \theta))$};
            \fill [myblue](2,1) circle (2pt);
            \fill [myblue](-1,2) circle (2pt);
            \node at (0, 2.5)[above right]{$\Im(z)$};
            \node at (2.5, 0)[above right]{$\Re(z)$};
            \node at (0,2)[left]{$y$};
            \node at (2,0)[below]{$x$};
            \node at (0,0)[below left]{$O$};
        \end{tikzpicture}
        \caption{Rotazione di \(z\) di \(90^\circ\) in senso antiorario del piano di Gauss}
        \label{fig:rotazione}
    \end{figure}
\end{curiosity}

\subsection{Notazione esponenziale}
\noindent\begin{minipage}{\textwidth}
    \begin{definition}[Notazione esponenziale]
        Un numero complesso \( z \) può essere scritto nella \textbf{forma esponenziale} come
        \[ z = re^{i\theta} = r(\cos\theta + i\sin\theta) \]
        \noindent Osservazioni e Proprietà:
        \begin{itemize}
            \item \(e^{i\theta} = \cos\theta + i\sin\theta\) è detta \textbf{formula di Eulero} \\
            \item \(e^{i\pi} + 1 = 0\) è detta \textbf{identità di Eulero}
        \end{itemize}
    \end{definition}
\end{minipage}

\begin{figure}[H]
    \centering
    \begin{tikzpicture}[scale=2]
        \pgfmathsetmacro{\rTheta}{180}
        \small
        \pgfmathsetmacro{\cosTheta}{cos(\rTheta)}
        \pgfmathsetmacro{\sinTheta}{sin(\rTheta)}
        \draw[->] (-1.5,0) -- (1.5,0) node[right] {$\Re$};
        \draw[->] (0,-1.5) -- (0,1.5) node[above] {$\Im$};
        \draw[-latex, ultra thick, myblue] (0,0) -- ({cos(\rTheta)}, {sin(\rTheta)});
        \draw[dashed] ({cos(\rTheta)}, {sin(\rTheta)}) -- ({cos(\rTheta)}, 0);
        \draw[dashed] ({cos(\rTheta)}, {sin(\rTheta)}) -- (0, {sin(\rTheta)});
        \draw(0,0) circle (1);
        \draw[-latex, myblue] (1.2,0) arc (0:180:1.2);
        \filldraw[black] ({cos(\rTheta)}, {sin(\rTheta)}) circle (1pt) node[anchor=south west] {$e^{i\rTheta}$};
        \node at (0, -2){$e^{i\rTheta} = \cos\rTheta + i\sin\rTheta = \cosTheta + i(\sinTheta)$};
        \node at (0, -2.5){$ exp(i\rTheta) = \sum_{n=0}^{+\infty} \frac{(i\rTheta)^n}{n!} = 1 + i\rTheta + \frac{(i\rTheta)^2}{2!} + ... + \frac{(i\rTheta)^n}{n!}$};
        \node at (0, -3){$\theta = \rTheta^\circ$};
    \end{tikzpicture}
    \caption{Identità di Eulero}
\end{figure}

\noindent\begin{minipage}{\textwidth}
    \begin{curiosity}[che significa \(e^{i\theta} = \cos\theta + i\sin\theta\)?]
        La formula di Eulero è una delle più belle formule della matematica, in quanto mette in relazione cinque numeri fondamentali della matematica: \(0\), \(1\), \(\pi\), \(i\) e \(e\).
        \begin{align*}
            e^{i\pi} + 1 & = 0
        \end{align*}
        Per tentare di capirla prima cosa sia \(e^x\) non si tratta di \(e \cdot e \cdot... \cdot e\) \(x\) volte, ma di
        \[e^{x} = \exp(x) = \sum_{n=0}^{+\infty} \frac{x^n}{n!} = 1 + x + \frac{x^2}{2!} + ... + \frac{x^n}{n!}\]
        che se \(x = i\theta\) diventa
        \[e^{i\theta} = \exp(i\theta) = \sum_{n=0}^{+\infty} \frac{(i\theta)^n}{n!} = 1 + i\theta + \frac{(i\theta)^2}{2!} + ... + \frac{(i\theta)^n}{n!}\]
        nella pratica se noi pnessimo \(\theta = 1\) risulterebbe
        \[e^{i} = \exp(i) = \sum_{n=0}^{+\infty} \frac{i^n}{n!} = 1 + i + \frac{i^2}{2!} + ... + \frac{i^n}{n!}\]
        ma cosa vuol dire elevare \(i^n\)? semplicemente come abbiamo visto prima risulta in un a rotazione di \(90^\circ\) in senso antiorario del piano di Gauss. visualemnte si può vedere come.
        \begin{figure}[H]
            \centering
            \begin{animateinline}[controls,loop]{60}
                \multiframe{30}{rTheta=0+12}{
                    \begin{tikzpicture}[scale=2]
                        \small
                        \pgfmathsetmacro{\cosTheta}{cos(\rTheta)}
                        \pgfmathsetmacro{\sinTheta}{sin(\rTheta)}
                        \draw[->] (-1.5,0) -- (1.5,0) node[right] {$\Re$};
                        \draw[->] (0,-1.5) -- (0,1.5) node[above] {$\Im$};
                        \draw[-latex, ultra thick, myblue] (0,0) -- ({cos(\rTheta)}, {sin(\rTheta)});
                        \draw[dashed] ({cos(\rTheta)}, {sin(\rTheta)}) -- ({cos(\rTheta)}, 0);
                        \draw[dashed] ({cos(\rTheta)}, {sin(\rTheta)}) -- (0, {sin(\rTheta)});
                        \draw(0,0) circle (1);
                        \draw[-latex, myblue] (1.2,0) arc (0:\rTheta:1.2);
                        \filldraw[black] ({cos(\rTheta)}, {sin(\rTheta)}) circle (1pt) node[anchor=south west] {$e^{i\rTheta}$};
                        \node at (0, -2){$e^{i\rTheta} = \cos\rTheta + i\sin\rTheta = \cosTheta + i(\sinTheta)$};
                        \node at (0, -2.5){$ exp(i\rTheta) = \sum_{n=0}^{+\infty} \frac{(i\rTheta)^n}{n!} = 1 + i\rTheta + \frac{(i\rTheta)^2}{2!} + ... + \frac{(i\rTheta)^n}{n!}$};
                        \node at (0, -3){$\theta = \rTheta^\circ$};
                    \end{tikzpicture}
                }
            \end{animateinline}
            \caption{Rotazione di \(e^{i\theta}\) di \(\theta\) in senso antiorario del piano di Gauss}
        \end{figure}
    \end{curiosity}
\end{minipage}
\begin{theorem}[Prodotto di due numeri complessi]
    Dati due numeri complessi \( z_1 = r_1 e^{i\theta_1} \) e \( z_2 = r_2 e^{i\theta_2} \), il loro \textbf{prodotto} è il numero complesso
    \[ z_1 \cdot z_2 = r_1 r_2 e^{i(\theta_1 + \theta_2)}. \]
\end{theorem}
\begin{theorem}[Quoziente di due numeri complessi]
    Dati due numeri complessi \( z_1 = r_1 e^{i\theta_1} \) e \( z_2 = r_2 e^{i\theta_2} \), il loro \textbf{quoziente} è il numero complesso
    \[ \frac{z_1}{z_2} = \frac{r_1}{r_2} e^{i(\theta_1 - \theta_2)}. \]
\end{theorem}

\begin{theorem}[Potenza di un numero complesso]
    Dato un numero complesso \( z = r e^{i\theta} \), la sua \textbf{potenza} è il numero complesso
    \[ z^n = r^n e^{i n \theta}. \]
\end{theorem}
\begin{theorem}[Radice di un numero complesso]
    Dato un numero complesso \( z = r e^{i\theta} \), la sua \textbf{radice} è il numero complesso
    \[ \sqrt[n]{z} = \sqrt[n]{r} e^{i \frac{\theta + 2k\pi}{n}}, \]
    dove \( k = 0, 1, 2, ..., n - 1 \).
\end{theorem}
\begin{theorem}[Formula di De Moivre]
    Dato un numero complesso \( z = r e^{i\theta} \), la sua \textbf{potenza} è il numero complesso
    \[ z^n = r^n e^{i n \theta}. \]
\end{theorem}
\begin{definition}[Esponenziale di un qualunque numero complesso]
    Dato un numero complesso \( z = x + iy \), la sua \textbf{esponenziale} è il numero complesso
    \[ e^z = e^{x + iy} = e^x e^{iy} = e^x (\cos y + i \sin y). \]
\end{definition}
\end{document}
